% s-182-p182-m001; original PDF 182, 183
\Aphorism{s-182-p182-m001}{271}
Auf zwei Arten kann ich an der \Sperr{Zeit} selbst — abgesehen vom erfüllenden Zeitinhalt — einen Genuss haben. Die eine Art, ich möchte sie die englisch-amerikanische nennen, erprobe ich in Grossstädten, wo es die Ausnutzung der Zeit bis in ihre kleinsten Theilchen gilt. Jenes aufreibende, nervenanspannende Hin und Her zwischen Post, Bahnhof, Dampfboot, das Aufeinanderbeziehen und Ineinanderdisponiren von Viertelstunden und Minuten, sodass keine müssig verrinnt, das Zusammendrängen eines Maximum von Arbeit in ein Minimum von Zeit — das ist, ehe es zur beständigen Plage wird, immerhin ein vorübergehendes Vergnügen. Die Muskeln spielen in einer ökonomischen Rhythmik, die jauchzen macht, und man hat das vollkommene Freiheitsgefühl der Improvisation — so ungefähr stelle ich mir die \Foreign{voluptas} des Seiltänzers oder des Bergsteigers an einer bösen Stelle vor. Fehlloses Ergreifen des Richtigen, ohne Wahl und Überlegung, ein allergeschwindester Kabelverkehr zwischen Gehirn und Organen, Motiv, Entschluss und That in Einem Augenblick: diese Zeitersparniss selbst ist Glück, ganz abgesehen vom Vortheil der ersparten Zeit. — Die andere Art Zeitgenuss, die italienische, lernt man im Süden; es ist die umgekehrte Kunst des Amerikaners, nämlich ein Minimum von Arbeit vertheilt auf ein Maximum von Zeit. Nichtsthun, nicht aus Mangel an Beschäftigung, nicht mit bösem Gewissen, sondern die bewusste Verlängerung des müssigen Augenblicks zur Zeitstrecke; Schlendern, Orangenessen, stundenlanges Liegen im Kahn, eine feierliche Zwecklosigkeit gleich dem Branden und Zurücklaufen der Meereswellen. Das südliche \Foreign{far }\Foreign{\Sperr{niente}} ist durch diese negative Benennung schlecht gezeichnet; es ist gar keine Negation, sondern das Positivste von der Welt — ich würde eher die Arbeit des Italieners ein Nichtfaullenzen nennen. Er ringt seinem Tagewerk, dem Nichtsthun, eine halbe Stunde Arbeit ab; das ist seine Ausnahme, sein zeitweiliges Heraustreten aus dem „ Milieu “, seine Selbstentäusserung, dabei fühlt er sich deplacirt, wie wir Nordländer uns umgekehrt bei unserem Nichtsthun deplacirt fühlen. Wir sagen vom Italiener mit einer gewaltsamen Metapher, er schlage die Zeit todt; aber ich finde, diesmal sind wir die Thierquäler, die mit Geissel und Stachel (mit „Beschäftigung“) das trottende Grauthier zu schnellerem Schritt spornen.
\par

